\section{Variational contacts and uniform phase confinement}\label{sec:contacts}

For the remainder of the paper, each array in the sequence is a global
maximizer of $C_{N_n}$ at its number of summands $N_n$ and satisfies
$L_n\to0$. Since the number of nonzero summands is at least $L_n^{-2}\to\infty$,
we have $N_n\to\infty$, and comparing with an $N_n$-summand i.i.d.\ Esseen
array gives $C_{N_n}\to\ce$, so Theorem~\ref{thm:near} applies.

\begin{lemma}[Variational contact and support coercivity]\label{lem:coercivity}
Normalize total variance to $B=1$, so $R=L$. For all sufficiently late global
maximizers, every support point of every summand satisfies
$|u|\le246R^{1/3}$. At an attaining positive closed threshold $z$, with
$D=F(z)-\Phi(z)$, $\alpha=\phi(z)$, $A_0=(3D+z\alpha)/2$, $v=\E X_i^2$,
$\rho=\E|X_i|^3$, and $M_i=\E[X_i|X_i|]$, every summand $X_i$ satisfies
\begin{equation}\label{eq:full-contact-unit}
 F_{-i}(z-u)+(\alpha+3CM_i)u+A_0u^2-C|u|^3
       \le F(z)+A_0v-C\rho\qquad(u\in\R),
\end{equation}
with equality at every $u\in\supp X_i$.
\end{lemma}
\begin{proof}
Perturbing $\Law(X_i)$ to $(1-e)\Law(X_i)+e\delta_u$ at the fixed closed
threshold $z$ and recentering and restandardizing gives \eqref{eq:full-contact-unit}
by maximality (adapting \citet[Section~2.1, eqs.~(7)--(12)]{HeChengCurrent2026}
to a single coordinate), and upper semicontinuity of $F_{-i}(z-u)$ extends
equality from a set of full $\Law(X_i)$-measure to $\supp X_i$. Substituting
the deleted Berry--Esseen bound $|F_{-i}(z-u)-\Phi((z-u)/\sqrt{1-v})|\le(123/10)R$
\citep[Theorem~2.1]{ChenShao2001} into the equality at $u=R^{1/3}w\in\supp X_i$
gives $(99/500)w^3\le1023/40+(129/10)w+(203/20)w^2$, which forces $w<246$.
\end{proof}

In the gap units of Lemma~\ref{lem:first-alias}, write $Y_i=J_i+Z_i$ with
$V=\sum_i\E Y_i^2$, $d=V^{-1/2}$, $P=\sum_i\E|Y_i|^3$, $Q=\sum_i\E Z_i^2$,
$T=\sum_i Z_i$, and standardized threshold $z=y/\sqrt V$. By
Theorem~\ref{thm:near} and Lemma~\ref{lem:coercivity}, $|Y_i|\le H=O(V^{1/3})$,
$\sup_i\E|Y_i|^3=o(V)$, and
\begin{equation}\label{eq:contact-inputs}
 \frac PV\to\re,\qquad \frac\kappa V\to\te,\qquad Q\to0,\qquad
 z\to0,\qquad \frac Dd\to\phz\left(\frac12+\frac\te6\right),\qquad
 \Cov(K,T)=o(\sqrt V).
\end{equation}

\begin{lemma}[Uniform support-phase confinement]\label{lem:whole-phase}
For the decomposition $Y_i=J_i+Z_i$,
\begin{equation}\label{eq:whole-phase}
       \eta_n:=\max_{1\le i\le N_n}\|Z_i\|_\infty\longrightarrow0,
\end{equation}
and the fractional part $\{y+\mu\}$ converges to $0$ modulo $1$. In
particular, whenever $t:=\E|T|>0$, we have $Q=o(t)$.
\end{lemma}
\begin{proof}
For any deleted coordinate $i$, put $v_i=\E Y_i^2$, $\rho_i=\E|Y_i|^3$,
$w_i=V-v_i$, and $\gamma_i=\kappa(Y_{-i})/w_i$, where $\kappa(Y_{-i})$ is the
deleted third-moment sum. For $|u|\le H$, if the deleted fractional part
$\tau_i(u)=\{y-u+\mu_{-i}\}$ lies in $[\epsilon,1-\epsilon]$, then
Lemma~\ref{lem:integer} applied to the deleted sum gives
\begin{equation}\label{eq:away-phase}
               F_{-i}(y-u)=G_i(u)-d\phz\tau_i(u)+o(d),
\end{equation}
where $G_i(u)=\Phi(x_i)+w_i^{-1/2}\phi(x_i)\{1/2+\gamma_i(1-x_i^2)/6\}$
($x_i=(y-u)/\sqrt{w_i}$) is the smooth deleted Edgeworth approximation. In
gap units, the contact score and level from \eqref{eq:full-contact-unit} are
\[
 g_i(u)=F_{-i}(y-u)+b_iu+A_0u^2-\lambda|u|^3,\qquad
 \ell_i=F(y)+A_0v_i-\lambda\rho_i,
\]
where
\[
 b_i=d\phi(z)+3\lambda\E[Y_i|Y_i|],\qquad
 \lambda=\frac DP=Cd^3,\qquad
 A_0=\frac{3D+z\phi(z)}{2V},
\]
so that $g_i(u)\le\ell_i$ for all $u\in\R$ with equality on $\supp Y_i$.
Expanding around $u=0$ shows that uniformly in $i$ and $|u|\le H$,
\begin{equation}\label{eq:smooth-contact-bound}
       G_i(u)+b_iu+A_0u^2-\lambda|u|^3-\ell_i\le o(d)
\end{equation}
(since $C\to\ce>\phz/6$ makes the cubic coefficient $-\lambda+\phz d^3/6$
strictly negative). If $u\in\supp Y_i$ had $u+\E K_i$ bounded away from $\Z$,
then $\tau_i(u)$ would stay in $[\epsilon,1-\epsilon]$ (since $\{y+\mu\}\to0$)
and \eqref{eq:away-phase}--\eqref{eq:smooth-contact-bound} would give
$g_i(u)<\ell_i$, contradicting $u\in\supp Y_i$.
Thus $\eta_n=\max_i\|Z_i\|_\infty\to0$. Finally, $\E T^4\le3Q^2+\eta_n^2Q$
and $t=\E|T|\ge Q^{3/2}/(\E T^4)^{1/2}\ge Q/\sqrt{3Q+\eta_n^2}$ imply
$Q=o(t)$ whenever $t>0$.
\end{proof}

\section{Relative conditional first-moment bounds}\label{sec:relative}

To compare the global maximizer against its exact integer rounding $J_i$ when
$t=\E|T|>0$ is arbitrarily small, we first adjust the common scale by
$1+\delta=1+o(A^{-1/2})$ to place the full characteristic product at its exact
local maximum near $2\pi$, keeping the integer labels $K_i$ unchanged.

\begin{lemma}[Alignment at the exact full-product peak]\label{lem:full-peak}
After a common rescaling by $1+o(A^{-1/2})$, the summands $Y_i=J_i+Z_i$
retain the same integer parts $J_i$ and satisfy \eqref{eq:whole-phase} and
\begin{equation}\label{eq:relative-hypotheses}
\begin{gathered}
 A=\Var K\to\infty,\qquad |J_i|\le H=O(A^{1/3}),\\
 \sum_{i=1}^N\E|J_i|^3=O(A),\qquad Q\to0,\qquad
                 |\Cov(K,T)|\le KHQ.
\end{gathered}
\end{equation}
\end{lemma}
\begin{proof}
Because $|\E e^{\ii(2\pi+x/\sqrt A)\sum_i J_i}|\to e^{-x^2/2}$ and $T\to0$ in
$L^1$, the full product $|\E e^{\ii s\sum_i Y_i}|$ attains a local maximum at
$s=2\pi(1+\delta)$ with $\delta=o(A^{-1/2})$. Rescaling $Y_i$ by $1+\delta$
and realigning the phases $\beta_i=(2\pi)^{-1}\arg\E e^{2\pi\ii Z_i^{\rm new}}$
keeps the nearest-integer labels $K_i$ unchanged, while stationarity
$\sum_i\E[Y_i^{\rm new}\sin(2\pi E_i)]/r_i=0$ at the full peak yields
$|\Cov(\sum_i Y_i^{\rm new},T_{\rm new})|\le KHQ_{\rm new}$, proving
\eqref{eq:relative-hypotheses}.
\end{proof}

\begin{lemma}[Multiplier bound in terms of $t=\E|\sum_i Z_i|$]\label{lem:multiplier}
Let $Z_1,\ldots,Z_N$ be independent centered random variables with
$t=\E|\sum_{i=1}^N Z_i|$. For any random variables $W_1,\ldots,W_N$ such that
each pair $(W_i,Z_i)$ is independent of $(W_j,Z_j)_{j\ne i}$ and
$\|W_i\|_\infty\le M_W$,
\begin{equation}\label{eq:multiplier}
 \left|\sum_{i=1}^N\E[W_iZ_i]\right|
       \le8t\left(\sum_{i=1}^N\E W_i^2\right)^{1/2}+5M_Wt.
\end{equation}
\end{lemma}
\begin{proof}
Let $Z_i'$ be an independent copy of $Z_i$ (independent of $W_i$) and set
$D_i=Z_i-Z_i'$. By Khintchine's inequality, $\E(\sum_i D_i^2)^{1/2}\le2\sqrt3\,t=:C_0t$.
Truncating at $b=4C_0t$, Markov's inequality gives
$\Pp(\max_{j\ne i}|D_j|\le b)\ge3/4$, so summing over the disjoint events
where exactly one $|D_i|>b$ yields $\sum_i\E[|D_i|\ind_{\{|D_i|>b\}}]\le(4/3)C_0t<5t$.
For $U=\sum_i\E[D_i^2\ind_{\{|D_i|\le b\}}]$, second/fourth-moment
interpolation gives $U/\sqrt{U+b^2}\le C_0t$, so $U<60t^2$. Applying
$\E[W_iZ_i]=\E[W_iD_i]$, Cauchy--Schwarz on $\{|D_i|\le b\}$, and
$\|W_i\|_\infty\le M_W$ on $\{|D_i|>b\}$ proves \eqref{eq:multiplier}.
\end{proof}

\begin{lemma}[Relative conditional first-moment bounds]\label{lem:conditional-relative}
Under the hypotheses of Lemma~\ref{lem:full-peak}, whenever $t=\E|T|>0$, for
each fixed $B_0>0$ and uniformly over $|k-\mu|\le B_0\sqrt A$,
\begin{equation}\label{eq:conditional-relative}
 c\,t\le\E(|T|\mid K=k)\le Kt,\qquad
                      |\E(T\mid K=k)|=o(t).
\end{equation}
\end{lemma}
\begin{proof}
Take two independent copies $(K,T)$ and $(K',T')$ of the row, with
$D_i=Z_i-Z_i'$ and $D=\sum_i D_i$. Fix a large constant $B>1$, truncate at
$b=Bt$, and set $d_i=D_i\ind_{\{|D_i|\le b\}}$, $W=\sum_i d_i$,
$U=\sum_i\E d_i^2$, and $\Delta=\sum_i\E[|D_i|\ind_{\{|D_i|>b\}}]$. As in the
proof of Lemma~\ref{lem:multiplier},
\begin{equation}\label{eq:two-copy-truncation}
 \Delta\le Kt,\qquad \sum_{i=1}^N\Pp(|D_i|>b)\le \frac KB,\qquad
                      U\le K_Bt^2\quad\bigl(K_B=O(B)\bigr).
\end{equation}
On the frequency cell $\|s\|\le c/H$ ($s=(s_1,s_2)$), set
\[
 f_i(u)=\E e^{\ii uJ_i},\qquad
 m_{ir}(s)=\frac{\E[d_i^re^{\ii(s_1J_i+s_2J_i')}]}{f_i(s_1)f_i(s_2)}.
\]
Applying
Lemma~\ref{lem:multiplier} to the antisymmetric variables $d_i$ bounds the
linear twisted sum by
\begin{equation}\label{eq:relative-Fourier}
 \left|\sum_{i=1}^N m_{i1}(s)\right|\le Kt\Bigl\{H\|s\|
  +(\sqrt{HA}+H^2)\|s\|^2+(H^2\sqrt A+H^3)\|s\|^3\Bigr\},
\end{equation}
which is $o(t)$ on the Gaussian scale $\|s\|=O(A^{-1/2})$. Fourier inversion
of the second and fourth moments of $W$ (using \eqref{eq:torus-damping}
outside $\|s\|\le c/H$) yields, uniformly for $|k_1-\mu|,|k_2-\mu|\le B_0\sqrt A$,
\begin{equation}\label{eq:two-count-moments}
 \E(W^2\mid K=k_1,K'=k_2)=U+o(t^2),\qquad
 \E(W^4\mid K=k_1,K'=k_2)\le C_Bt^4.
\end{equation}
For the discarded tail $D-W=\sum_i D_i\ind_{\{|D_i|>b\}}$, one-pair deletion
ratios $\Pp(K_{-i}=k-k_{i1},K'_{-i}=k-k_{i2})/\Pp(K=K'=k)=1+o(1)$ uniformly
over shifts $O(H)$ imply $\E(|D-W|\mid K=K'=k)\le2\Delta\le2Kt$, while
\eqref{eq:two-count-moments} and $U\le K_Bt^2$ give
$\E(|W|\mid K=K'=k)\le K(B)t$, so $\E(|D|\mid K=K'=k)\le K'(B)t$. For the
lower bound, if $\Delta\ge\epsilon t$, conditioning on the disjoint events
where a single $|D_i|>b$ gives $\E(|D|\mid K=K'=k)\ge c\Delta\ge c\epsilon t$
for large $B$; if $\Delta<\epsilon t$, then $U\ge(1-\epsilon)^2t^2$ and
moment interpolation in \eqref{eq:two-count-moments} gives
$\E(|W|\mid K=K'=k)\ge c_Bt$, so choosing $\epsilon<c_B/4$ gives
$2\E(|T|\mid K=k)\ge\E(|D|\mid K=K'=k)\ge c_Bt/2$. Finally, applying the
one-copy Fourier expansion with \eqref{eq:relative-Fourier} and
$|\Cov(K,T)|\le KHQ=o(Ht)$ gives $|\E(T\mid K=k)|=o(t)$, and conditional
Jensen yields
\[
 \E(|T|\mid K=k)\le\E(|T-T'|\mid K=K'=k)+|\E(T\mid K=k)|\le K''(B)t.
\]
\end{proof}

\section{Exact phase collapse for global maximizers}\label{sec:comparison}

For the central integer layer $k=\operatorname{round}(y+\mu)$ and offset
$\theta=y+\mu-k\to0$, define
\[
 w(\theta)=\Pp(T>\theta\mid K=k),\qquad
                        f(\theta)=\theta+w(\theta).
\]

\begin{lemma}[Relative closed-threshold loss]\label{lem:closed-loss}
If $t=\E|T|>0$, then at the attaining central layer $k$ and offset $\theta$,
\begin{equation}\label{eq:relative-closed-loss}
        f(\theta)\ge ct,\qquad
                   |\theta|+w(\theta)+t\le K f(\theta).
\end{equation}
\end{lemma}
\begin{proof}
Truncating $T$ to $X=T\ind_{\{|T|\le1/4\}}$ changes its conditional CDF,
mean, and $L^1$ norm by $O(Q)=o(t)$ (via \eqref{eq:all-layer}), so
$m=\E(X\mid K=k)=o(t)$ and $s=\E(|X|\mid K=k)\ge ct$ by
Lemma~\ref{lem:conditional-relative}. Put $u=\Pp(X>\theta\mid K=k)$. For
$0\le\theta\le1/4$,
$s=2\E(X_+\mid K=k)-m\le2\theta+\frac12u-m\le2(\theta+u)-m$;
for $\theta=-b\in[-1/4,0)$, $m\le(1/4+b)u-b\le u/2-b$ and $s\le u/2-m$ give
$\theta+u\ge s+2m$, proving \eqref{eq:relative-closed-loss}.
\end{proof}

\begin{proposition}[Exact collapse of rounding residuals]\label{prop:collapse}
For every sufficiently late global maximizer in the sequence, $T=0$ almost
surely, so every summand $Y_i=J_i$ is supported on a single coset of $\Z$.
\end{proposition}
\begin{proof}
Suppose $t=\E|T|>0$ along a subsequence. Compare the global maximizer
$(Y_1,\ldots,Y_N)$ directly against the centered integer array
$(J_1,\ldots,J_N)$ at the same number of summands $N$ and integer threshold
$k-\mu$. By \eqref{eq:all-layer}, decomposing $\{S\le k+\theta-\mu\}$ by the
value of $K$ gives
\begin{equation}\label{eq:exact-crossing}
 F_Y(k+\theta-\mu)=F_K(k)-p_kw(\theta)+E_{\rm cross},\qquad
                         |E_{\rm cross}|\le KdQ=o\bigl(df(\theta)\bigr).
\end{equation}
Since $p_k=d\phi(z)+o(d)$, $\Phi((k+\theta-\mu)/\sqrt V)=\Phi((k-\mu)/\sqrt A)+d\phi(z)\theta+o(df(\theta))$
(using $|V-A|\le KHQ=o(At)$), and Lemma~\ref{lem:closed-loss} controls both
$|\theta|$ and $w(\theta)$ by $f(\theta)$, the discrepancies satisfy
\begin{equation}\label{eq:selected-comparison}
          D_Y=D_J-d\phi(z)f(\theta)+o\bigl(df(\theta)\bigr).
\end{equation}
Along the interpolation $Y_i(s)=J_i+sZ_i$ ($0\le s\le1$),
Lemma~\ref{lem:multiplier} with $W_i=3Y_i(s)|Y_i(s)|$ bounds
$\bigl|\frac{d}{ds}\sum_i\E|Y_i(s)|^3\bigr|\le K(\sqrt{HA}+H^2)t=o(At)$, so
$|P_Y-P_J|=o(At)$. Since $V\sim A$ ($d=V^{-1/2}$) and $P_Y\sim P_J\sim\re V$,
the prefactors $V^{3/2}/P_Y\sim(\re d)^{-1}$ and $A^{3/2}/P_J$ differ by a
relative error of $o(t)$, that is, $|V^{3/2}/P_Y-A^{3/2}/P_J|=o(d^{-1}t)$.
Because $D_J=O(d)$ and $t\le Kf(\theta)$, multiplying
\eqref{eq:selected-comparison} by $V^{3/2}/P_Y$ shows that
$(J_1,\ldots,J_N)$ has Berry--Esseen ratio strictly larger than that of
$(Y_1,\ldots,Y_N)$ by at least $c'f(\theta)\ge c''t>0$, contradicting global
maximality of $(Y_1,\ldots,Y_N)$. Thus $t=0$ eventually, so $Z_i=0$ almost
surely for all $i$.
\end{proof}

\section{Discrete concavity and reduction to two adjacent lattice points}\label{sec:lattice}

By Proposition~\ref{prop:collapse}, for all sufficiently late global
maximizers, each summand $Y_i$ is supported on a single coset of $\Z$.

\begin{lemma}[Reduction to two adjacent lattice points]\label{lem:adjacent}
For all sufficiently late global maximizers, every nonconstant summand $Y_i$
is supported on two adjacent points of its lattice coset.
\end{lemma}
\begin{proof}
Let $c_i(x)=\Pp(\sum_{j\ne i}Y_j=x)$ be the deleted point mass function on
its coset, with deleted variance $w_i=V-v_i$, $d_i=w_i^{-1/2}$, deleted
signed cubic sum $\kappa_i=\kappa(Y_{-i})$, and $H_3(x)=x^3-3x$. Fourier
inversion with multiplier $e^{-\ii u}-1$ and \eqref{eq:torus-damping} gives,
uniformly in $i$ and $x$,
\begin{equation}\label{eq:pmf-first-difference}
 c_i(x+1)-c_i(x)=d_i^2\phi'(d_ix)+\frac{d_i^3}2\phi''(d_ix)
       +d_i^3\frac{\kappa_i}{6w_i}(H_3\phi)'(d_ix)+o(d^3).
\end{equation}
For the contact function $g_i(u)=F_{-i}(y-u)+b_iu+A_0u^2-\lambda|u|^3$ on the
coset of $Y_i$, its centered second difference
$\Delta^2g_i(u)=g_i(u+1)-2g_i(u)+g_i(u-1)$ has deleted-CDF contribution
$c_i(y-u+1)-c_i(y-u)$, whose leading $-d^2z\phi(z)$ term cancels the
$d^2z\phi(z)$ part of $2A_0$, leaving for $|u|\le H+1$:
\begin{equation}\label{eq:lattice-concavity}
 d^{-3}\Delta^2g_i(u)=\phz(1+u)-\ce
       \bigl\{|u+1|^3+|u-1|^3-2|u|^3\bigr\}
       +o(1+|u|)+O(d|z|v_i).
\end{equation}
Because $\ce>\phz$ and $|u+1|^3+|u-1|^3-2|u|^3$ equals $6|u|$ for $|u|\ge1$
and $2+6u^2-2|u|^3\ge2+4u^2$ for $|u|\le1$, the leading expression in
\eqref{eq:lattice-concavity} is bounded above by $-c(1+|u|)$ with $c>0$.
Applying the discrete Green identity on any interval of length $L\ge2$
between two contact points $u_0,u_0+L\in\supp Y_i$,
\[
 g_i(u_0+j)-g_i(u_0)
 =-\sum_{k=1}^{L-1}\frac{\min(j,k)(L-\max(j,k))}{L}\Delta^2g_i(u_0+k)\le0,
\]
first bounds $L$ and $v_i$ by a uniform constant (making $d|z|v_i=o(1)$ in
\eqref{eq:lattice-concavity}), and then implies $\Delta^2g_i(u)<0$ strictly
at every coset point, which makes $L\ge2$ impossible. Thus $L=1$ for every
nonconstant summand.
\end{proof}

Since every nonconstant summand is supported on two points at distance $1$,
all active summands share a common support diameter. Applying Theorem~1.1 of
\citet{BellBernoulli2026} at the same horizon $N_n$ shows that all active
summands are identically distributed, completing the proof of
Theorem~\ref{thm:global}.
